\chapter{Discussion}
\label{ch:discussion}

This chapter steps back from individual figures and tables to ask what the
experiments in Chapters~\ref{ch:results} and~\ref{ch:scalability} actually
established, where they fell short of the paper's own results, and why.

\section{What worked}

\textbf{The efficiency claim holds up, in a setting the paper never
tested.} The paper's argument for eALS is that it avoids the $O(K^3)$
matrix inversion of classical whole-data ALS while staying exactly as
parallelizable. Both halves were checked directly: Section~\ref{sec:scale_k}
measured $10.9\times$ growth in training time for a $16\times$ growth in
$K$, nowhere close to the $4096\times$ a genuinely cubic solver would need.
Section~\ref{sec:parallel}'s independence argument, that the parallel and
sequential update rules are mathematically identical and not merely
``close enough'', was validated against a synthetic dataset
(Section~\ref{sec:correctness}) with no special-casing for the distributed
execution path.

\medskip
\textbf{Exact parallelism is a real advantage over gradient-based
alternatives.} Every worker in the user-update sweep reads the same frozen
$Q$ and $S^q$ and writes to a disjoint slice of $P$, so there is no
analogue of the gradient-conflict problem that makes naively parallelized
SGD either slow (if synchronized) or approximate (if not). This is the
paper's own argument, but building the actual PySpark version confirmed
that it survives contact with a real distributed runtime rather than
remaining a property of the equations alone.

\medskip
\textbf{The optimal missing-data weight $c_0=512$ transferred exactly.}
Given how differently this project's dataset was curated compared to the
paper's 2016 snapshot (a third of the users, a fifth of the interactions,
see Table~\ref{tab:dataset}), it was not obvious that an absolute-magnitude
hyperparameter like $c_0$ would land on the same optimum. That it did
(Figure~\ref{fig:weighting_c0}) is a reassuring sign that the underlying
mechanism, balancing observed-data loss against a much larger pool of
weighted missing entries, behaves consistently across dataset scale, at
least along this one axis.

\section{What did not work as well, and why}
\label{sec:weaknesses}

\textbf{Popularity-aware weighting's advantage over uniform weighting was
much smaller here than in the paper.} The paper reports a statistically
significant improvement from $\alpha=0$ to $\alpha\approx0.4$, confirmed
with a paired $t$-test across repeated trials; this project found only a
marginal, statistically untested difference (Section~\ref{sec:weighting}).
Three factors plausibly contribute, though this project's setup cannot
fully separate them: the dataset is roughly a fifth the size of the
paper's, and popularity skew tends to be less extreme in smaller,
more heavily-filtered samples; the sampled-negative evaluation protocol
(Chapter~\ref{ch:protocol}) is a coarser ranking signal than the paper's
full-catalog ranking and could plausibly wash out a small real effect; and
each configuration here was trained once, with no repeated runs to average
out initialization noise, unlike the paper's repeated-trial significance
tests. Repeating every configuration and re-running full-catalog evaluation
at scale were not practical within this project's resources, but both
would be the direct way to determine whether the gap is a property of the
dataset or of these measurement choices.

\medskip
\textbf{Parallelism stopped helping past four partitions, on this
machine.} Section~\ref{sec:scale_par} found the sweet spot at 4 partitions
on a 4-core laptop, with 8 and 16 performing slightly worse. This is a
finding about the hardware, not the algorithm: the update rule stays
exactly as parallelizable at 16 partitions as at 4
(Section~\ref{sec:parallel}'s argument does not depend on partition
count), but Spark's per-task scheduling overhead does not shrink with the
work each task has to do, so past a point adding partitions only adds
overhead. On a real multi-node cluster the constraint shifts from
per-task scheduling to inter-node transfer of the broadcast $Q$/$S^q$
matrices, so this experiment's specific numbers would not transfer, though
the underlying diminishing-returns trade-off almost certainly would.
Confirming that on real cluster hardware was outside what this project
could test.

\medskip
\textbf{The inner update loop is plain Python, not vectorized across
users.} Listing~\ref{lst:update_users} loops over one user at a time
inside each partition, even though a single user's update is only a
handful of NumPy operations over a short neighbor list. That per-user
Python overhead is the most likely reason MLlib's JVM-native ALS is
$3$-$6\times$ faster in absolute terms across every $K$ tested in
Section~\ref{sec:scale_mllib}, despite implementing a costlier theoretical
algorithm. Batching several users' updates into fewer, larger NumPy calls
would likely close much of that gap without touching the parallelization
strategy or the underlying math at all. This was not attempted here, since
closing an implementation-efficiency gap is a different question from
validating the paper's algorithmic claims, which is what this project
prioritized.

\medskip
\textbf{Scope: offline eALS only, not the paper's online extension.} As
flagged in Chapter~\ref{ch:intro}, this project implements the paper's
Algorithm~1 and its parallelization, not Algorithm~2 (the incremental
online update) or the online evaluation protocol of the paper's
Section~5.3. This matches the assignment brief, which asks specifically
for an efficient, parallelizable MF-ALS algorithm, but it does mean this
report cannot speak to the ``online recommendation'' half of the paper's
title.

\medskip
Two of these gaps, the smaller popularity-weighting effect and the
non-matching absolute accuracy numbers, trace back to the same root
cause: a differently-curated, roughly five-times-smaller dataset than the
paper's, combined with a coarser sampled-negative evaluation used to keep
the scalability study tractable. Both choices are stated where they are
made (Chapter~\ref{ch:protocol}), and their combined effect is that this
report's accuracy numbers are internally consistent with each other but
should not be read as digit-for-digit comparable to the paper's own Table
and Figure values.
